% ============================================================
%  CHAPTER 8 — RESULTS AND DISCUSSION
% ============================================================
\chapter{Results}


\section{Emotion Recognition Results}

\subsection{Model Training Convergence}

The ResNet-18 model was trained for 30 epochs on the RAF-DB dataset. The
training loss converged steadily, and the validation accuracy stabilised after
approximately 20 epochs. The best model checkpoint was achieved at epoch 24,
recording the highest validation accuracy before early stopping halted training.

\begin{table}[H]
\centering
\caption{Training and Validation Metrics at Key Epochs}
\begin{tabular}{cccc}
\toprule
\textbf{Epoch} & \textbf{Training Loss} & \textbf{Val.\ Accuracy (\%)} & \textbf{Val.\ Loss} \\
\midrule
5              & 1.42                   & 71.3                         & 1.18 \\
10             & 0.98                   & 78.6                         & 0.89 \\
15             & 0.74                   & 82.1                         & 0.74 \\
20             & 0.58                   & 84.3                         & 0.68 \\
24 (best)      & 0.48                   & 85.7                         & 0.63 \\
30             & 0.43                   & 85.4                         & 0.65 \\
\bottomrule
\end{tabular}
\end{table}

\subsection{Per-Class Classification Performance}

The confusion matrix analysis revealed that the model performs best on
\textit{Happy} (93\% precision) and \textit{Neutral} (88\% precision),
which are the most frequently occurring classes in the RAF-DB dataset.
\textit{Disgust} and \textit{Fear} showed relatively lower performance
due to class imbalance, despite the use of weighted sampling.

\begin{table}[H]
\centering
\caption{Per-Class F1 Scores on RAF-DB Test Set}
\begin{tabular}{lccc}
\toprule
\textbf{Emotion} & \textbf{Precision} & \textbf{Recall} & \textbf{F1-Score} \\
\midrule
Angry            & 0.81               & 0.78            & 0.79 \\
Disgust          & 0.74               & 0.70            & 0.72 \\
Fear             & 0.76               & 0.73            & 0.74 \\
Happy            & 0.93               & 0.95            & 0.94 \\
Neutral          & 0.88               & 0.89            & 0.88 \\
Sad              & 0.82               & 0.80            & 0.81 \\
Surprise         & 0.87               & 0.86            & 0.86 \\
\midrule
\textbf{Weighted Avg.} & \textbf{0.86} & \textbf{0.857} & \textbf{0.857} \\
\bottomrule
\end{tabular}
\end{table}

\subsection{Comparison with Existing Approaches}

The performance of the MindTrace ResNet-18 model is compared against
representative methods from the literature reviewed in Chapter 2:

\begin{table}[H]
\centering
\caption{Comparison of Emotion Recognition Accuracy on RAF-DB}
\begin{tabular}{lll}
\toprule
\textbf{Method}                  & \textbf{Architecture}   & \textbf{Accuracy (\%)} \\
\midrule
Li et al.\ (2017)                & DLP-CNN                 & 84.13 \\
Wang et al.\ (2020)              & SCN                     & 87.03 \\
Zhao et al.\ (2021)              & TransFER (Transformer)  & 90.91 \\
\textbf{MindTrace (Ours)}        & ResNet-18 fine-tuned    & \textbf{85.70} \\
\bottomrule
\end{tabular}
\end{table}

The MindTrace model achieves competitive accuracy compared to the DLP-CNN
baseline and closely approaches the SCN method, while remaining significantly
more lightweight and suitable for real-time inference. The Transformer-based
TransFER achieves higher accuracy but is computationally expensive and
unsuitable for low-latency deployment on constrained hardware.

\section{Engagement Scoring Results}

\subsection{Score Behaviour Under Controlled Conditions}

The engagement scoring algorithm was validated by having participants perform
predefined actions while the system was running. The observed engagement scores
matched expected behaviour across all tested conditions:

\begin{table}[H]
\centering
\caption{Engagement Score Observations Under Controlled Conditions}
\begin{tabular}{lc}
\toprule
\textbf{Participant Action}                     & \textbf{Observed Score Range} \\
\midrule
Looking directly at screen (neutral expression) & 78 -- 85 \\
Smiling / happy expression, facing camera       & 88 -- 96 \\
Looking away (yaw $>$ 25°)                      & 30 -- 45 \\
Eyes closed / frequent blinking                 & 40 -- 55 \\
Visible frustration / angry expression          & 55 -- 65 \\
Combined distraction (look away + tired)        & 20 -- 35 \\
\bottomrule
\end{tabular}
\end{table}

The results confirm that the composite score correctly differentiates
between engaged and disengaged states. The highest scores correspond to
direct eye contact with a positive emotional state, while combined
distraction cues (poor head pose and high blink rate) produce the lowest
scores, as expected.

\subsection{Real-Time Session Timeline}

A typical 10-minute session produces a continuous timeline of engagement
scores. Analysis of five recorded sessions showed:
\begin{itemize}
    \item Average session engagement: \textbf{72.4 out of 100}
    \item Most frequent emotion detected: \textbf{Neutral} (43\% of frames)
    \item Second most frequent: \textbf{Happy} (28\% of frames)
    \item Distraction events (score $<$ 40): \textbf{6.2\%} of frames
          on average
\end{itemize}

\section{System Performance Results}

\subsection{API Response Latency}

End-to-end latency was measured from the moment the frontend sends a frame
to when the response is rendered on the dashboard:

\begin{table}[H]
\centering
\caption{End-to-End Latency Breakdown}
\begin{tabular}{lc}
\toprule
\textbf{Stage}                          & \textbf{Average Time (ms)} \\
\midrule
Frame capture and Base64 encoding       & 12 \\
Network transmission (client to server) & 38 \\
Face detection (MediaPipe)              & 28 \\
Emotion model inference (ResNet-18)     & 95 \\
Engagement computation                  & 14 \\
Network transmission (server to client) & 31 \\
\midrule
\textbf{Total End-to-End Latency}       & \textbf{218 ms} \\
\bottomrule
\end{tabular}
\end{table}

The total latency of 218 ms comfortably satisfies the SRS requirement of
$<$ 300 ms per analysis cycle and supports a real-time update rate of
approximately 4--5 frames per second on the dashboard.

\subsection{System Scalability}

The stateless backend design was stress-tested by simulating concurrent
users using the Locust load-testing framework:

\begin{table}[H]
\centering
\caption{Scalability Test Results}
\begin{tabular}{ccc}
\toprule
\textbf{Concurrent Users} & \textbf{Mean Response Time (ms)} & \textbf{Error Rate (\%)} \\
\midrule
1                         & 187                              & 0.0 \\
5                         & 204                              & 0.0 \\
10                        & 243                              & 0.0 \\
20                        & 298                              & 0.8 \\
50                        & 512                              & 4.2 \\
\bottomrule
\end{tabular}
\end{table}

The system handles up to 20 concurrent users within acceptable latency bounds.
Beyond 20 users, response times increase and occasional timeouts occur due to
the single-instance deployment on Hugging Face Spaces. Horizontal scaling
(multiple instances behind a load balancer) would address this limitation.

\section{System User Interface}

This section showcases the final user interface of the MindTrace application,
demonstrating the user journey from authentication to real-time analytics.

\subsection{Authentication and Onboarding}

The system provides a modern, responsive interface for user onboarding.
Figure \ref{fig:onboarding} shows the landing page and the secure login/signup
modules.

\begin{figure}[H]
\centering
\begin{subfigure}[b]{0.48\textwidth}
    \centering
    \includegraphics[width=\textwidth]{ss_landing}
    \caption{MindTrace Landing Page}
\end{subfigure}
\hfill
\begin{subfigure}[b]{0.48\textwidth}
    \centering
    \includegraphics[width=\textwidth]{ss_login}
    \caption{Secure Login Interface}
\end{subfigure}
\caption{System Onboarding and Authentication}
\label{fig:onboarding}
\end{figure}

\subsection{Live Monitoring Dashboard}

The core functionality of MindTrace is demonstrated through the live dashboard.
Figure \ref{fig:live_dashboard} captures the system performing real-time
emotion inference and focus tracking via the student's webcam.

\begin{figure}[H]
\centering
\includegraphics[width=0.85\textwidth]{ss_live_monitoring}
\caption{Real-Time Monitoring with Emotion Overlay}
\label{fig:live_dashboard}
\end{figure}

\subsection{Session Analytics and History}

After each session, the system provides detailed behavioural insights.
Figure \ref{fig:analytics} shows the post-session dashboard featuring emotion
distribution and engagement trends.

\begin{figure}[H]
\centering
\includegraphics[width=0.85\textwidth]{ss_analytics}
\caption{Post-Session Behavioural Analytics Dashboard}
\label{fig:analytics}
\end{figure}

\section{Cloud Deployment and Database}

MindTrace is deployed as a distributed cloud application to ensure global
accessibility and reliable data persistence.

\begin{figure}[H]
\centering
\begin{subfigure}[b]{0.48\textwidth}
    \centering
    \includegraphics[width=\textwidth]{ss_vercel}
    \caption{Frontend Hosting on Vercel}
\end{subfigure}
\hfill
\begin{subfigure}[b]{0.48\textwidth}
    \centering
    \includegraphics[width=\textwidth]{ss_huggingface}
    \caption{Backend API on Hugging Face}
\end{subfigure}
\caption{Cloud Deployment Infrastructure}
\label{fig:deployment}
\end{figure}

The session data is persisted in a MongoDB Atlas cluster, as shown in Figure
\ref{fig:mongodb}, which allows for real-time data exploration and long-term
storage of student engagement records.

\begin{figure}[H]
\centering
\includegraphics[width=0.85\textwidth]{ss_mongodb}
\caption{MongoDB Atlas Data Explorer (Session Records)}
\label{fig:mongodb}
\end{figure}

\section{Discussion}

\subsection{Achievement of Objectives}

All primary objectives defined in Chapter 1 were successfully achieved:

\begin{itemize}
    \item \textbf{Emotion classification}: Seven-class emotion detection
          with 85.7\% accuracy was implemented and deployed in real time.
    \item \textbf{Head pose and blink detection}: Accurately flags
          distraction events with a low false-positive rate.
    \item \textbf{Engagement scoring}: The composite score provides
          meaningful and interpretable feedback that correlates with
          observed participant behaviour.
    \item \textbf{Data persistence}: Session records are reliably stored
          and retrievable from MongoDB Atlas.
    \item \textbf{Dashboard}: A fully functional React dashboard with
          live overlays, charts, and admin panel was delivered.
    \item \textbf{Security}: JWT authentication and bcrypt password
          hashing were implemented and validated.
    \item \textbf{Cloud deployment}: The system is live and accessible
          via public URLs on Hugging Face Spaces and Vercel.
\end{itemize}

\subsection{Limitations}

\begin{itemize}
    \item The model exhibits reduced accuracy in low-lighting conditions,
          as the RAF-DB training images are primarily well-lit.
    \item \textit{Disgust} and \textit{Fear} classes have lower F1 scores
          due to underrepresentation in the dataset.
    \item The engagement score weights ($w_e$, $w_p$, $w_b$) were set
          empirically and may not generalise equally across all demographic
          groups.
    \item The single-instance deployment limits concurrent user capacity
          to approximately 20 users.
\end{itemize}
